\paragraph{Transport to the fixed complementary twist.}
Put $d_k=|k|$, so that $k=-d_k<0$. The old source has sign $-1$;
its Fourier coefficient at the positive index $d$ tests $E^{-d}$.
Choose the new source for $E^k$ with sign $+1$, normalize its
coefficient at $1$ to be $1$, and impose a support containing the
old support and all primes dividing $k$. For the following comparison
$D$ is positive and squarefree, is prime to that support, and lies
in the square unit classes of $1$ required by the new filter.
In particular $D$ is a square in every old support completion, so
$d_kD$ and $d_k$ lie in the same old filter. Waldspurger comparison
on the two sources gives
\[
 \begin{split}
  v_2 a_{\rm new}(D)
    &=v_2 a_{\rm old}(d_kD)-v_2 a_{\rm old}(d_k)\\
    &\ge \omega(D).
 \end{split}
\]
The equality is a valuation equality; the two sources may have
different root-of-unity phases. The normalized central-value ratios
agree, since both are
\[
 \sqrt D\,\frac{L(E^{kD},1)}{L(E^k,1)}.
\]
Fixed period constants cancel in each ratio. The old unit condition
at $d_k$, and $\omega(d_kD)=\omega(d_k)+\omega(D)$, prove the
inequality. The signed twist is $kD$ throughout; no Fourier
coefficient at a negative index is used. The final family already
satisfies these unit classes. In the nonempty-seed argument the
fresh prime $p'$ is chosen to match all fixed labels of $p$, so
$pp'$ satisfies the same filter, including its dyadic condition.

For a unit of this normalized source equality holds, so
$v_2L_*(h_xk)\le2\omega(h_x)+C_{e,\rm final}$ with one fixed
period-normalization constant.
Its $K$-split seed must be nonempty. Suppose otherwise. The integral
coefficient valuations then improve the preceding lower bound to
$\omega(D)+1$ on every nonempty product of allowed $K$-split primes.
For a fresh good $K$-split prime $p$ outside the support, apply the
